\section{The object of study: a transformer in equations}
\label{sec:transformer}

Before we can interpret a model we need a precise description of what it computes. This section defines a
decoder-only transformer of the GPT-2 type \citep{vaswani2017attention,radford2019language} and extracts the
structural facts on which the methods of the later sections rest. The presentation follows the ``residual
stream'' view of \citet{elhage2021mathematical}.

\subsection{From text to a prediction}

A \emph{tokenizer} cuts the input text into pieces called \emph{tokens}, which are words or fragments of words,
and replaces each piece by its index in a fixed vocabulary of size $V$. The input of the model is therefore a
sequence of indices $t_1,\dots,t_n$ with $t_i\in\{1,\dots,V\}$. For every position $i$ the model returns a vector
of $V$ real numbers called \emph{logits}. We write $\logit_i(t)$ for the logit of token $t$ at position $i$. The
predicted probability that the next token $t_{i+1}$ equals $t$ is the softmax of the logits,
\begin{equation}\label{eq:softmax}
\Pr(t_{i+1}=t \mid t_1,\dots,t_i) \;=\; \frac{\exp(\logit_i(t))}{\sum_{t'=1}^{V}\exp(\logit_i(t'))} .
\end{equation}
Training adjusts the weights so as to minimise the average of $-\log\Pr(t_{i+1}\mid t_1,\dots,t_i)$ over a large
corpus of text. This quantity is called the cross-entropy loss, and it is small when the true next token receives a
high probability.

\begin{exercise}[\extitle{Softmax sees only differences}{derivation}]\label{ex:softmax}
(a) Show that adding the same number to all logits at position $i$ does not change the probabilities
\eqref{eq:softmax}. (b) Show that the ratio of the probabilities of two tokens $t$ and $t'$ equals
$\exp(\logit_i(t)-\logit_i(t'))$. (c) Let $t^\star$ be the true next token. Show that the derivative of the loss
$-\log\Pr(t_{i+1}=t^\star\mid t_1,\dots,t_i)$ with respect to $\logit_i(t)$ equals
$\Pr(t_{i+1}=t\mid t_1,\dots,t_i)-1$ if $t=t^\star$ and $\Pr(t_{i+1}=t\mid t_1,\dots,t_i)$ if $t\ne t^\star$.
(d) What do (a) to (c) suggest about the right way to measure how well a model separates two candidate tokens?
\end{exercise}

\subsection{The residual stream}

The model keeps, for every position $i$, a vector in $\R^d$ that is updated layer by layer. This vector is called
the \emph{residual stream} at position $i$. We write $x_i^{(\ell)}\in\R^d$ for its value at the input of layer
$\ell$, where the layers are numbered $\ell=0,\dots,L-1$, and $x_i^{(L)}$ for its final value.
\Cref{fig:residual} shows the flow of the computation.

\paragraph{Embedding.} The stream starts as the sum of a vector for the token and a vector for the position,
\begin{equation}\label{eq:embed}
x_i^{(0)} \;=\; W_E[t_i] + W_{\mathrm{pos}}[i],
\end{equation}
where $W_E[t]\in\R^d$ denotes the column of the embedding matrix $W_E\in\R^{d\times V}$ that belongs to token $t$,
and $W_{\mathrm{pos}}[i]\in\R^d$ is a learned vector for position $i$.

\paragraph{Layer normalisation.} Every component reads the stream through the map
\begin{equation}\label{eq:ln}
\LN(x) \;=\; \frac{x-\mean(x)\,\one}{\scale(x)},\qquad
\mean(x)=\frac1d\,\one^\top x,\qquad
\scale(x)=\frac{1}{\sqrt d}\,\big\|x-\mean(x)\,\one\big\|,
\end{equation}
where $\one\in\R^d$ is the vector of ones and $\|\cdot\|$ is the Euclidean norm. $\LN$ subtracts the average coordinate and rescales the result to a
fixed length. The original definition also multiplies by a learned vector and adds another one. Since $\LN$ is
always followed by a linear map, these two learned vectors can be absorbed into that map and its bias, and we do
so throughout.

\paragraph{Attention heads.} Layer $\ell$ contains $H$ attention heads, numbered $h=0,\dots,H-1$. A head is the
only kind of component that moves information between positions. Head $h$ of layer $\ell$ has four weight matrices
$W_Q,W_K,W_V\in\R^{d_{\mathrm{head}}\times d}$ and $W_O\in\R^{d\times d_{\mathrm{head}}}$ (each head has its own,
and we suppress the indices $\ell,h$ on them). The width $d_{\mathrm{head}}$ of a head is much smaller than
$d$. They enter only through the two products
\begin{equation}\label{eq:qkov}
W_{QK} \;=\; W_Q^\top W_K\in\R^{d\times d},\qquad W_{OV} \;=\; W_O W_V\in\R^{d\times d}.
\end{equation}
For a position $i$ and every earlier or equal position $j\le i$, the head computes a score, turns the scores into
weights by a softmax over $j$, and outputs a weighted sum:
\begin{align}
A^{\ell,h}_{ij} &\;=\; \frac{\exp\!\big(\LN(x_i^{(\ell)})^\top W_{QK}\,\LN(x_j^{(\ell)})/\sqrt{d_{\mathrm{head}}}\big)}
{\sum_{j'\le i}\exp\!\big(\LN(x_i^{(\ell)})^\top W_{QK}\,\LN(x_{j'}^{(\ell)})/\sqrt{d_{\mathrm{head}}}\big)},
\label{eq:pattern}\\
\head_i^{\ell,h} &\;=\; \sum_{j\le i}A^{\ell,h}_{ij}\;W_{OV}\,\LN(x_j^{(\ell)}) \;\in\R^d. \label{eq:headout}
\end{align}
The numbers $A^{\ell,h}_{ij}\ge 0$ sum to one over $j$ and are called the \emph{attention pattern} of the head.
We call $i$ the \emph{destination} and $j$ the \emph{source}. In the usual terminology,
$W_Q\LN(x_i^{(\ell)})$ is the \emph{query} of the destination, $W_K\LN(x_j^{(\ell)})$ is the \emph{key} of the source
and $W_V\LN(x_j^{(\ell)})$ is its \emph{value}. The score is the inner product of query and key, so a source
receives a large weight if its key matches what the destination asks for. The two products in \eqref{eq:qkov} have separate
jobs. $W_{QK}$ decides \emph{from where} position $i$ reads, and $W_{OV}$ decides \emph{what} is written to
position $i$ once a source has been chosen. The restriction $j\le i$ ensures that a position never sees later
tokens.

\paragraph{MLP.} After the heads, layer $\ell$ applies a two-layer perceptron to every position separately. With
$\tilde x_i^{(\ell)}$ denoting the stream after the heads of layer $\ell$,
\begin{equation}\label{eq:mlp}
\mlp_i^{\ell} \;=\; W_{\mathrm{out}}\,\mathrm{GELU}\big(W_{\mathrm{in}}\,\LN(\tilde x_i^{(\ell)})+b_{\mathrm{in}}\big)+b_{\mathrm{out}} \;\in\R^d,
\end{equation}
with weight matrices $W_{\mathrm{in}}$, $W_{\mathrm{out}}$ and bias vectors $b_{\mathrm{in}}$, $b_{\mathrm{out}}$ that
belong to the layer. Here $\mathrm{GELU}$ is a smooth variant of $\max(0,\cdot)$ applied to each coordinate. The coordinates of the
vector inside $\mathrm{GELU}$ are the \emph{neurons} of the layer.

\paragraph{Update rule and output.} Both kinds of components \emph{add} their output to the stream,
\begin{equation}\label{eq:update}
\tilde x_i^{(\ell)} \;=\; x_i^{(\ell)}+\sum_{h=0}^{H-1}\head_i^{\ell,h},\qquad
x_i^{(\ell+1)} \;=\; \tilde x_i^{(\ell)}+\mlp_i^{\ell},\qquad \ell=0,\dots,L-1 .
\end{equation}
The logits are a linear function of the normalised final stream. With $W_U[t]\in\R^d$ denoting the row of the
unembedding matrix $W_U\in\R^{V\times d}$ that belongs to token $t$, written as a column vector,
\begin{equation}\label{eq:unembed}
\logit_i(t) \;=\; W_U[t]^\top\,\LN\big(x_i^{(L)}\big).
\end{equation}
Two kinds of bias terms are omitted from these formulas. The linear maps inside a head have bias vectors, and
absorbing the learned shift of the final $\LN$ into the unembedding leaves a constant $b_U[t]$ that is added to the
logit of token $t$. All statements below stay true with them, if the output bias of each attention layer is counted
as one more component and $b_U[t]$ as one more summand of $\logit_i(t)$.

\begin{figure}[t]
\centering
\begin{tikzpicture}[>=Latex, font=\small,
  box/.style={draw=tutgrey, rounded corners=2pt, minimum height=6.5mm, inner sep=4pt, fill=#1},
  box/.default=white, add/.style={circle, draw=tutgrey, inner sep=0.5pt, fill=white}]
\draw[line width=1.8pt, tutblue] (0,0.35) -- (0,5.55);
\node[box=tutpaper] (emb) at (0,0) {embedding $W_E[t_i]+W_{\mathrm{pos}}[i]$};
\node (tok) at (0,-0.95) {token $t_i$};
\draw[->, tutgrey] (tok) -- (emb);
\node[add] (p1) at (0,2.35) {$+$};
\node[add] (p2) at (0,4.35) {$+$};
\node[box=tutblue!12] (att) at (4.2,1.55) {$\LN$, then the $H$ attention heads};
\node[box=tutred!12] (mlp) at (3.25,3.55) {$\LN$, then the MLP};
\draw[->, tutgrey] (0,1.55) -- (att.west);
\draw[->, tutgrey] (att.north) |- (p1);
\draw[->, tutgrey] (0,3.55) -- (mlp.west);
\draw[->, tutgrey] (mlp.north) |- (p2);
\node[align=left] (others) at (9.3,1.55) {streams $x_j^{(\ell)}$ of the\\ positions $j\le i$};
\draw[->, tutgrey] (others.west) -- (att.east);
\node[box=tutpaper] (un) at (0,5.9) {$\LN$, then unembedding $W_U$};
\node (out) at (0,6.85) {logits for the token after position $i$};
\draw[->, tutgrey] (un) -- (out);
\node[left] at (-0.1,1.0) {$x_i^{(\ell)}$};
\node[left] at (-0.1,3.0) {$\tilde x_i^{(\ell)}$};
\node[left] at (-0.1,5.0) {$x_i^{(\ell+1)}$};
\draw[decorate, decoration={brace, amplitude=5pt}, tutgrey] (-1.5,0.75) -- (-1.5,5.2)
  node[midway, left=6pt, align=right] {layer $\ell$,\\ repeated for\\ $\ell=0,\dots,L-1$};
\end{tikzpicture}
\caption{The residual stream at one position $i$ (thick blue line). Each component reads the stream through
$\LN$ and adds its output back. Only the attention heads see the streams of other positions.}
\label{fig:residual}
\end{figure}

\paragraph{Sizes.} GPT-2 small has $L=12$ layers, $H=12$ heads per layer, stream width $d=768$,
$d_{\mathrm{head}}=64$, $3072$ neurons per MLP, a vocabulary of $V=50257$ tokens and about $124$ million
parameters. It is customary to name a head by layer and index, so that ``head 9.6'' is head $h=6$ of layer
$\ell=9$. Current large models have the same architecture up to details and are larger by two to three orders of
magnitude.

\begin{exercise}[\extitle{Attention by hand}{derivation}]\label{ex:byhand}
Consider one head with $d=2$ and $d_{\mathrm{head}}=1$ in a model without $\LN$ and without biases. Its weights are
$W_Q=(0\ \ 2)$, $W_K=(1\ \ 0)$, $W_V=(1\ \ 0)$ and $W_O=(0\ \ 1)^\top$. The stream at the input of the layer is
$x_1^{(\ell)}=(1,0)^\top$ at position $1$ and $x_2^{(\ell)}=(0,1)^\top$ at position $2$. (a) Compute $W_{QK}$ and
$W_{OV}$. (b) Compute the attention pattern and the output of the head at the destination $i=2$. (c) Do the same
at the destination $i=1$. (d) The attention weights of a destination always sum to one, so a head cannot decide
to read from nowhere. How can a model still arrange that a head writes almost nothing at some destinations?
\end{exercise}

\begin{exercise}[\extitle{Counting parameters}{derivation}]\label{ex:params}
Count the weights of one attention head of GPT-2 small, without biases. Using the total of $124$ million
parameters, which fraction of the model is taken by the token embedding matrix $W_E$?
\end{exercise}

\begin{exercise}[\extitle{What one head can write}{derivation}]\label{ex:rank}
(a) Show that $W_{OV}$ has rank at most $d_{\mathrm{head}}$. What does this say about the set of vectors that a
single head of GPT-2 small can write to the stream? (b) In GPT-2 small, $d_{\mathrm{head}}=d/H$. What is the
largest possible dimension of the space to which the $H$ heads of one layer can write together?
\end{exercise}

\subsection{Four structural facts}
\label{sec:facts}

The following facts are immediate from the definitions, and each of them is the basis of a family of methods.

\begin{proposition}[The final stream is a sum of component outputs]\label{prop:sum}
For every position $i$,
\begin{equation}\label{eq:sum}
x_i^{(L)} \;=\; W_E[t_i]+W_{\mathrm{pos}}[i]+\sum_{\ell=0}^{L-1}\Big(\sum_{h=0}^{H-1}\head_i^{\ell,h}+\mlp_i^{\ell}\Big).
\end{equation}
\end{proposition}
\begin{proof}
We write the final stream as a telescoping sum over the layers.
\begin{align*}
x_i^{(L)} &= x_i^{(0)}+\sum_{\ell=0}^{L-1}\big(x_i^{(\ell+1)}-x_i^{(\ell)}\big) && \by{telescoping sum}\\
&= x_i^{(0)}+\sum_{\ell=0}^{L-1}\Big(\sum_{h=0}^{H-1}\head_i^{\ell,h}+\mlp_i^{\ell}\Big) && \by{by \eqref{eq:update}}\\
&= W_E[t_i]+W_{\mathrm{pos}}[i]+\sum_{\ell=0}^{L-1}\Big(\sum_{h=0}^{H-1}\head_i^{\ell,h}+\mlp_i^{\ell}\Big). && \by{by \eqref{eq:embed}}\qedhere
\end{align*}
\end{proof}

\begin{proposition}[Each logit is a sum over components, given the final scale]\label{prop:dla}
Let $\mathcal G_i$ be the set of the $2+L(H+1)$ vectors on the right-hand side of \eqref{eq:sum}. We call them
the \emph{components} at position $i$. Then for every token $t$,
\begin{equation}\label{eq:dla}
\logit_i(t) \;=\; \sum_{g\in\mathcal G_i}\frac{W_U[t]^\top\big(g-\mean(g)\,\one\big)}{\scale\big(x_i^{(L)}\big)} .
\end{equation}
\end{proposition}
\begin{proof}
We insert the definition of $\LN$ and the sum of \cref{prop:sum} into the read-out.
\begin{align*}
\logit_i(t) &= W_U[t]^\top\,\LN\big(x_i^{(L)}\big) && \by{by \eqref{eq:unembed}}\\
&= \frac{W_U[t]^\top\big(x_i^{(L)}-\mean(x_i^{(L)})\,\one\big)}{\scale(x_i^{(L)})} && \by{by \eqref{eq:ln}}\\
&= \frac{W_U[t]^\top\big(\sum_{g\in\mathcal G_i}g-\mean\big(\sum_{g\in\mathcal G_i}g\big)\,\one\big)}{\scale(x_i^{(L)})} && \by{by \cref{prop:sum}}\\
&= \sum_{g\in\mathcal G_i}\frac{W_U[t]^\top\big(g-\mean(g)\,\one\big)}{\scale(x_i^{(L)})}. && \by{since $\mean$ is linear}\qedhere
\end{align*}
\end{proof}

The summand of a component $g$ in \eqref{eq:dla} is called its \emph{direct logit attribution}. We use it in
\cref{sec:lens}. Two more facts need no proof.

\begin{itemize}
\item \textbf{Information moves between positions only through attention.} By \eqref{eq:mlp}, the output of an
MLP at position $i$ depends on the stream at position $i$ alone, and the same holds for $\LN$ and for the
unembedding. If the final position ``knows'' something about an earlier token, some attention head has carried
it there.
\item \textbf{The coordinates of the stream have no special meaning.} Every component reads the stream through
a linear map after $\LN$ and writes to it through a linear map. Let $\Omega$ be an orthogonal matrix with $\Omega\one=\one$.
Then $\LN(\Omega x)=\Omega\,\LN(x)$. If we replace the stream by $\Omega x$, multiply every reading matrix by
$\Omega^\top$ from the right and every writing matrix by $\Omega$ from the left, the model computes the same
function. Hence nothing singles out coordinate number
$17$ of the stream. The MLP neurons are different, because $\mathrm{GELU}$ acts on each coordinate separately.
We say that the neurons form a \emph{privileged basis} and that the residual stream has none.
\end{itemize}

\begin{intuition}
Think of the residual stream as a shared whiteboard with $d$ dimensions of space. Each head and each MLP reads
what is on the board, computes something, and adds its result to the board. The final answer is read off the
board by one linear map. \Cref{prop:sum,prop:dla} say that we can account for the final answer contribution by
contribution. The last fact says that the board has no natural grid: the meaningful ``places'' on it are
directions that we have to find.
\end{intuition}

\begin{exercise}[\extitle{A direction that nobody reads}{derivation}]\label{ex:ones}
Let $s\in\R$. Show that $\LN(x+s\,\one)=\LN(x)$. Conclude that replacing the stream $x_i^{(\ell)}$ by
$x_i^{(\ell)}+s\,\one$ at any single layer and position does not change any logit.
\end{exercise}

\begin{exercise}[\extitle{Rotating the stream}{derivation}]\label{ex:rotation}
(a) Prove the claim made above that $\LN(\Omega x)=\Omega\,\LN(x)$ for every orthogonal matrix $\Omega$ with
$\Omega\one=\one$. (b) Show by an example with $d=2$ that the claim fails for a general orthogonal matrix.
(c) A colleague plots coordinate number $17$ of the residual stream over many prompts and finds that it is large
on texts about sports. Use the last fact of this subsection to explain why this observation says little about
how the model represents sports.
\end{exercise}

\subsection{The running example in this notation}
\label{sec:metric}

In the IOI task of \cref{sec:running} we look at the last position $n$ and compare two tokens, the indirect object
(the correct name) and the subject (the repeated name). Our measure of performance is the \emph{logit difference}
\begin{equation}\label{eq:ld}
\logit_n(\mathrm{IO})-\logit_n(\mathrm{S}) \;=\; \big(W_U[\mathrm{IO}]-W_U[\mathrm{S}]\big)^\top\LN\big(x_n^{(L)}\big).
\end{equation}
By \eqref{eq:softmax}, the ratio of the two predicted probabilities equals the exponential of this difference.
On our 12 prompts the logit difference of GPT-2 small is $3.70$ on average and lies between $2.06$ and $5.53$.
A logit difference of $3.70$ on a prompt means that the model judges the correct name about
$e^{3.70}\approx 40$ times as likely as the repeated one on that prompt.

We prefer the logit difference to the probability of the correct name for two reasons. It is a linear function of
$\LN(x_n^{(L)})$, so \cref{prop:dla} applies to it term by term. And it removes everything that the two names have
in common. A component that only says ``the next token is a name'' raises both logits equally and cancels in
\eqref{eq:ld}.

\begin{exercise}[\extitle{What the logit difference hides}{concept}]\label{ex:metric}
Suppose that some head raises the logits of all first names by the same amount at the last position of every IOI
prompt. (a) What is its contribution to the logit difference \eqref{eq:ld}? (b) Is this head irrelevant for the
task? (c) Propose a metric under which its contribution becomes visible.
\end{exercise}

\subsection{Why the problem is hard}

Three properties of the object make interpretation difficult, and each later section responds to one of them.
The first is \emph{size}. GPT-2 small already has $144$ heads and $36{,}864$ MLP neurons, and each of them acts
at every position. The second is that information is \emph{distributed}. By the last fact of \cref{sec:facts},
a concept need not sit in one coordinate or one neuron (\cref{sec:features,sec:sae}). The third is
\emph{nonlinearity}. The attention pattern \eqref{eq:pattern}, the $\mathrm{GELU}$ and the scale in $\LN$ make
the effect of a component depend on what all the others are doing, which is why we will need interventions
(\cref{sec:interventions}) and cannot stop at decompositions.

\begin{takeaways}
A transformer keeps one vector per position, the residual stream, and every component adds to it. The final
stream is a sum of component outputs, and each logit is a sum of per-component terms once the final scale is
fixed. Attention heads are the only route between positions. The stream has no privileged coordinates, and the
MLP neurons do.
\end{takeaways}

\begin{readinglist}
\item \citet{vaswani2017attention}, \emph{NeurIPS}. The original description of the architecture. It predates the
residual-stream view, and comparing the two presentations is instructive.
\item \citet{weiss2021thinking}, \emph{ICML}. Defines a small programming language whose operations correspond to
attention and to MLP layers. Writing a task in this language shows how many layers and heads a transformer needs
for it.
\item \citet{geva2021transformer}, \emph{EMNLP}, and \citet{geva2022transformer}, \emph{EMNLP}. Two studies of the
MLP layers. The first reads their input weights as patterns that are detected and their output weights as tokens
that are promoted. The second describes the output of an MLP as a sum of updates to the stream, read through the
unembedding.
\item \citet{dar2023analyzing}, \emph{ACL}. Interprets the weights of a transformer directly, by expressing them
in the space of token embeddings, without running the model on any input.
\item \citet{lindner2023tracr}, \emph{NeurIPS}. A compiler that turns a known algorithm into the weights of a
transformer. The resulting models have a known ground truth, on which interpretability methods can be tested.
\item \citet{csordas2025language}, \emph{NeurIPS}. Uses the additive structure of the residual stream to ask how
much each layer of current large models contributes, and finds that the later layers contribute much less than
the earlier ones.
\end{readinglist}
